\documentclass[english, parskip=full]{scrartcl}
\usepackage[utf8]{inputenc}  % Kodierung der Datei
\usepackage[english]{babel}  % Sprache auf Deutsch 
\usepackage{color}
\usepackage{graphicx, gensymb}
\usepackage{amsfonts, cancel, amssymb}
\usepackage{amsmath, amsthm, array, esint}
\usepackage{booktabs, multirow}  % schönere Tabellen
\usepackage{caption}  % Anpassung der Captions
\usepackage{scrlayer-scrpage}    % Manipulation des Seitenstils
\pagestyle{scrheadings}  % Stil für die Seite setzen
\automark{section}  % Abschnittsnamen als Seitenbeschriftung 
\setheadsepline{.5pt}    % Kopzeile durch Linie abtrennen
\usepackage[hidelinks]{hyperref}  % Links und weitere PDF-Features
\usepackage{typearea, tikz, pgffor, verbatim}
\usetikzlibrary{calc, arrows.meta,arrows, graphs, quotes}
\usepackage[cache=false, outputdir=build]{minted}
\usemintedstyle{vs}
\usepackage{placeins} % provides \FloatBarrier

\definecolor{bg}{rgb}{0.9,0.9,0.9}
\newcommand\numberthis{\addtocounter{equation}{1}\tag{\theequation}}
\usepackage{svg}
\usepackage{siunitx}
\usepackage{physics}
\usepackage{tensor}
\renewcommand{\bra}{}
\newcommand{\kom}{}
\newcommand{\brak}{}
\renewcommand{\braket}{}
\newcommand{\intx}{}
\newcommand{\intr}{}
\newcommand{\kla}{}
\newcommand{\klam}{}
\newcommand{\klammer}{}
\newcommand{\oper}{}
\newcommand{\tecxt}{}
\newcommand{\Bra}{}
\newcommand{\Ket}{}
\renewcommand{\i}{\text{i}}
\newcommand{\e}{\text{e}}
\newcommand*{\Fourier}{\textsc{Fourier}\ }
\newcommand*{\F}{\mathcal{F}}
\newcommand*{\sinc}{\text{sinc\,}}
\newcommand{\conv}{\mathop{\scalebox{3.5}{\raisebox{-0.38ex}{\makebox[0.5\width][c]{$\ast$}}}}\limits}

\newcommand*{\titel}{Lagrangian for electrodynamics and applications}
\newcommand*{\autor}{Joachim Schwardt}

\hypersetup{pdfauthor={\autor}, pdftitle={\titel}}

%\titlehead{titlehead}
\subject{Theoretical Electrodynamics}
\title{\titel}
\author{\autor}
\date{\today}

\begin{document}

%\begin{titlepage}
\maketitle  % Titel setzen
%\tableofcontents  % Inhaltsverzeichnis setzen
%\end{titlepage}

\section{Lagrangian formulation of electrodynamics}
We want to start our presentation on electrodynamics by postulating the Lagrangian $\mathcal{L}_\tecxt\text{ED} = \mathcal{L}_0 + \mathcal{L}_\tecxt\text{int}$ where
\begin{align}
\mathcal{L}_0 &:= -\frac{1}{2} F^{\mu\nu}F_{\mu\nu},\\
\mathcal{L}_\tecxt\text{int} &:= \mu_0 J^\mu A_\mu
\end{align}
and $F^{\mu\nu} = \partial^\mu A^\nu - \partial^\nu A^\mu$.
Here $(A^\mu) = (\phi, \textbf{A})^\top$ is the four-vector potential, and $(J^\mu) = (\rho, \mathbf{j})^\top$ is the four-vector current density.
The continuity equation can be written as $\partial_\mu J^\mu =0$, and the Lorenz-gauge is given by $\partial_\mu A^\mu=0$.
Expanding the free Lagrangian leads to
\begin{align}
\mathcal{L}_0 &= -\frac{1}{2} \kla\left( \partial^\mu A^\nu - \partial^\nu A^\mu \right) \kla\left( \partial_\mu A_\nu - \partial_\nu A_\mu \right) = \\
&= -\frac{1}{2}\kla\big( \partial^\mu A^\nu\partial_\mu A_\nu - \partial^\nu A^\mu \partial_\mu A_\nu + (\mu\leftrightarrow\nu) \big) = \\
&= -A^\nu \partial^\mu \partial_\mu A_\nu - \kla\left( \partial^\mu A^\nu \right) \partial_\mu A_\nu + A^\mu \partial_\mu \underbrace{\partial^\nu A_\nu}_{=\,0} + \kla\left( \partial^\nu A^\mu  \right) \partial_\mu A_\nu = \\
&= -A^\nu \square A_\nu - \kla\left( \partial^\mu \phi \right) \partial_\mu \phi + \kla\left( \partial^\mu \textbf{A} \right) \cdot \partial_\mu \textbf{A} + \kla\left( \partial^\nu\phi \right)\partial_t A_\nu + \kla\left( \partial^\nu \textbf{A} \right) \cdot \nabla A_\nu = \\
&= -A^\nu \square A_\nu  - \kla\left( \partial^\mu \phi \right) \partial_\mu \phi + \kla\left( \partial^\mu \textbf{A} \right) \cdot  \partial_\mu \textbf{A} + (\partial_t\phi)^2 + 2(\nabla \phi)  \cdot \partial_t\textbf{A} + \nabla \cdot \kla\left( \textbf{A}\cdot\nabla \right) \textbf{A}.
\end{align}
We may now derive the free field equation for the electrostatic potential $\phi$ 
\begin{align}
0 &= \frac{\partial\mathcal{L}_0}{\partial\phi} - \partial_\mu \frac{\partial \mathcal{L}_0}{\partial \partial_\mu \phi} = -\square \phi + 2\partial_\mu \kla\left( \partial^\mu \phi \right) - \underbrace{2\partial_t \kla\left( \partial_t\phi \right) - 2\nabla \cdot \partial_t\textbf{A}}_{=\, 2\partial_t \partial_\mu A^\mu\,=\,0} =  \square \phi
\end{align}
and that for the vector potential $\textbf{A}$
\begin{align}
0 &= \frac{\partial\mathcal{L}_0}{\partial\textbf{A}} - \partial_\mu \frac{\partial \mathcal{L}_0}{\partial \partial_\mu \textbf{A}} = \square \textbf{A} - 2\partial_\mu \kla\left( \partial^\mu \textbf{A} \right) -2\partial_t (\nabla\phi) - \textbf{e}_j \partial_m \frac{\partial \kla\left( \partial_k A_l\partial_l A_k \right)}{\partial (\partial_m A_j)} = \\
&= -\square \textbf{A} - 2\nabla \kla\left( \partial_t\phi \right) - \textbf{e}_j \partial_m \kla\left( \partial_k A_l \delta_{ml}\delta_{jk} + \partial_l A_k \delta_{mk}\delta_{jl} \right) = \\
&= -\square \textbf{A} - 2\nabla (\partial_t\phi) - 2\textbf{e}_j\partial_j (\nabla\cdot \textbf{A}) = -\square \textbf{A}.
\end{align}
We may combine the two equations into $\square A^\mu = 0$.
Adding the interaction then simply yields
\begin{align}
\square A^\mu &= \mu_0 J^\mu. \label{eqn:four potential PDE}
\end{align}


\section{Green's function for the d'Alembert operator}
Since \autoref{eqn:four potential PDE} is an inhomogeneous equation, we can solve it using Green's method.
Assume that
\begin{align}
\square G(\textbf{r},t) &= \delta (\textbf{r})\delta (t). \label{eqn:dAlembert Green PDE}
\end{align}
Then, for a given charge density $J^\mu$, we may write the solution to \autoref{eqn:four potential PDE} as 
\begin{align}
A^\mu(\textbf{r},t) &= \mu_0 \intr\int_{\mathbb{R}^{4}} J^\mu(\textbf{r}',t') G(\textbf{r}-\textbf{r}', t-t')\, \mathrm{d}^{3}\textbf{r}'\,\tecxt\text{d}t'.
\end{align}
This is obvious from the definition of $G$ in \autoref{eqn:dAlembert Green PDE}.
\\
In order to compute the Green's function, we take the Fourier transformation of the PDE.
This yields
\begin{align}
1 &= \F\klam\left[ \square G(\textbf{r},t) \right](\textbf{k},\omega) =  \klam\left[ (\i\omega)^2 - (\i\textbf{k})^2 \right] G(\textbf{k},\omega),\\
\Rightarrow\ \ \ G(\textbf{k},\omega) &= \frac{1}{\textbf{k}^2 - \omega^2}.
\end{align}
Note that we use the same letter for $G(\textbf{r},t)$ and its Fourier transformation, but we avoid possible confusion by explicitly writing out the arguments.
We may now solve this equation by the inverse Fourier transformation
\begin{align}
G(\textbf{r},t) &= \frac{1}{(2\pi)^4} \intr\int_{\mathbb{R}^{4}} \frac{\e^{-\i (\omega t - \textbf{k}\cdot \textbf{r})} }{\textbf{k}^2 - \omega^2} \, \mathrm{d}^{3}\textbf{k}\, \tecxt\text{d}\omega = \\
&= \frac{1}{(2\pi)^4} \intr\int_{\mathbb{R}^{ }} \intx\int_{0}^{\infty} \intx\int_{0}^{2\pi} \intx\int_{0}^{\pi} \frac{\e^{-\i\omega t} \e^{\i kr\cos\theta}}{k^2 - \omega^2} \, \sin\theta\mathrm{d}\theta\, \mathrm{d}\phi\, k^2\mathrm{d}k \, \mathrm{d}^{ }\omega = \\
&= \frac{1}{(2\pi)^3} \intr\int_{\mathbb{R}^{ }} \e^{-\i\omega t} \intx\int_{0}^{\infty} \frac{k^2}{k^2 - \omega^2} \frac{\e^{\i kr \cos\theta}}{-\i kr}\bigg\vert_{\theta=0}^{\theta=\pi}\, \mathrm{d}k \, \mathrm{d}^{ }\omega = \\
&= \frac{-\i}{(2\pi)^3 r} \intr\int_{\mathbb{R}^{ }} \e^{-\i\omega t} \intx\int_{0}^{\infty} \frac{k}{k^2 - \omega^2} \kla\left( \e^{\i kr} - \e^{-\i kr} \right) \, \mathrm{d}k \, \mathrm{d}^{ }\omega = \\
&= \frac{-\i}{(2\pi)^3 r} \intr\int_{\mathbb{R}^{ }} \e^{-\i\omega t} \intr\int_{\mathbb{R}^{ }} \frac{k\e^{\i kr}}{k^2 - \omega^2} \, \mathrm{d}^{ }k \, \mathrm{d}^{ }\omega . \label{eqn:dAlembert Green partial result}
\end{align}
We may solve the $k$-integration with the residual theorem.
Let $\epsilon > 0$ be some small quantity, in the sense that we want to later take the limit $\epsilon \rightarrow 0$.
Define the integral
\begin{align}
I_\epsilon &:= \intr\int_{\mathbb{R}^{ }} \frac{k\e^{\i kr}}{k^2 - \omega^2 + \i\epsilon} \, \mathrm{d}^{ }k.
\end{align}
Note that the denominator may be factorized as $(k-k_+)(k-k_-)$, where $k_\pm = \pm\sqrt{\omega^2 - \i\epsilon}$.
Note that for $\epsilon> 0$ only $k_-$ has a positive imaginary part.
\autoref{figure:tikz contour for integration} shows the contour $\Gamma$ along which we want to integrate.
\begin{figure}[!ht]
\centering
\begin{tikzpicture}[scale=1]    
    \pgfmathsetmacro{\omega}{1.0}
    \pgfmathsetmacro{\eps}{0.3}
    \pgfmathsetmacro{\r}{0.03}
    \pgfmathsetmacro{\R}{3}
    
    
    \coordinate (o) at (0,0);
    \coordinate (xmax) at (5,0);
    \coordinate (xmin) at (-5,0);
    \coordinate (ymax) at (0,3.5);
    \coordinate (ymin) at (0,-1);
    
    \coordinate (kmax) at (\R,0);
    \coordinate (kmin) at (-\R,0);
    \coordinate (knorthwest) at (45:\R);
    
    \coordinate (k+) at (\omega,-\eps);
    \coordinate (k-) at (-\omega,\eps);
    
    \node[anchor=north] at (xmax) {Re\,$k$};
    \node[anchor=west] at (ymax) {Im\,$k$};
    
%    \draw[fill, color=black!35!white] (o) circle (\R);
%    \draw (o) circle (\R);
%    \node[below=0.5cm, right] at (0,\R) {$R$};
%    \draw[color=black, -{Stealth[length=1.5mm,width=1.5mm]}] (o) to (0,\R);
%    \node[below=0.5cm, left] at (o) {$M$};
%    
    \draw[fill] (k+) circle (\r);
    \node[below] at (k+) {$k_+$};
    
    \draw[fill] (k-) circle (\r);
    \node[above] at (k-) {$k_-$};
    
    \draw[thick, -{Stealth[length=1.5mm,width=1.5mm]}] (kmin) arc (180:30:\R);
    \draw[thick] (knorthwest) arc (45:0:\R);
    \node[right=0.3cm] at (knorthwest) {$\Gamma$};
    \draw[thick, -{Stealth[length=1.5mm,width=1.5mm]}] (kmax) to (o);
    \draw[thick] (o) to (kmin);
    
    \draw[thick, color=blue, -{Stealth[length=1.0mm,width=1.0mm]}] (o) to (knorthwest);
    \node[above=0.2cm, left=0.1cm, color=blue] at (45:0.5*\R) {$R$};
    
%    
%    \draw[color=black, thick=0.05cm, -{Stealth[length=1mm,width=1mm]}] (r0) .. controls (r1) and (r2)  .. (r3);
%    \draw[color=blue, thick=0.05cm, -{Stealth[length=1.5mm,width=1.5mm]}] (r0) to (r0right); 
%    \node[above, color=blue] at (r0right) {$\vec{v}_0$};      
%    
    
    \draw[color=black, -{Stealth[length=1.5mm,width=1.5mm]}] (xmin) to (xmax);
    \draw[color=black, -{Stealth[length=1.5mm,width=1.5mm]}] (ymin) to (ymax);
\end{tikzpicture}
\caption{For $\epsilon> 0$ only the pole at $k=k_-$ is inside the contour $\Gamma$.}
\label{figure:tikz contour for integration}
\end{figure}
From the residual theorem we have
\begin{align}
I_\Gamma &:= \intx\int_{\Gamma }^{ } \frac{k\e^{\i kr}}{(k-k_+)(k-k_-)} \, \mathrm{d}k = 2\pi\i \lim_{k\rightarrow k_-} \frac{k\e^{\i kr}}{k - k_+} = 2\pi\i \frac{k_-\e^{\i k_-r}}{k_- - k_+} \overset{\epsilon\,\rightarrow\,0}{\longrightarrow} \pi\i\e^{-\i\omega r}.
\end{align}
Define the integral along the arc as $I_R$ and substitute $k=R\e^{\i t}$ for $t\in (0,\pi)$
\begin{align}
I_R &:= \intx\int_{0}^{\pi} \frac{R\e^{\i t} \e^{\i rR\e^{\i t}}}{R^2\e^{2\i t} - \omega^2 + \i \epsilon}\, R\i\e^{\i t}\mathrm{d}t.
\end{align}
We want to show that this integral vanishes for $R\rightarrow\infty$.
Consider the absolute value
\begin{align}
|I_R| &\le \intx\int_{0}^{\pi} \frac{R^2\e^{-rR\sin(t)}}{|R^2\e^{2\i t} - \omega^2 + \i\epsilon |}\, \mathrm{d}t \le \intx\int_{0}^{\pi} \frac{R^2\e^{-rR\sin (t)}}{R^2 - |\omega^2 + \i\epsilon|}\, \mathrm{d}t = \\
&= \frac{2R^2}{R^2 - |\omega^2 + \i\epsilon|} \intx\int_{0}^{\pi/2} \e^{-rR\sin(t)}\, \mathrm{d}t \le \frac{2R^2}{R^2 - |\omega^2 + \i\epsilon|} \intx\int_{0}^{\pi/2}\e^{-rR \frac{2}{\pi} t}\, \mathrm{d}t = \\
&= \frac{2R^2}{R^2 - |\omega^2 + \i\epsilon|} \frac{\e^{-rR} - 1}{-rR\frac{2}{\pi}} \overset{R\,\rightarrow\,\infty}{\longrightarrow} 0.
\end{align}
Since $I_\Gamma = I_\epsilon + I_R$, we now find
\begin{align}
I_\epsilon &= I_\Gamma = \pi\i\e^{-\i \omega r}. \label{eqn:residual integral I eps}
\end{align}
One may easily repeat this calculation for $\epsilon < 0$, which merely amounts to a different sign in the exponent of  \autoref{eqn:residual integral I eps}.
\\
Inserting our results into \autoref{eqn:dAlembert Green partial result} leads to
\begin{align}
G(\textbf{r},t) &= \lim_{\epsilon\rightarrow 0} \frac{-\i}{(2\pi)^3 r} \intr\int_{\mathbb{R}^{ }} \e^{-\i\omega t} \intr\int_{\mathbb{R}^{ }} \frac{k\e^{\i kr}}{k^2 - \omega^2 \pm \i\epsilon} \, \mathrm{d}^{ }k \, \mathrm{d}^{ }\omega = \frac{\pi }{(2\pi)^3 r} \intr\int_{\mathbb{R}^{ }} \e^{-\i\omega t} \e^{\pm\i\omega r} \, \mathrm{d}^{ }\omega = \\
&= \frac{1}{4\pi r} \frac{1}{2\pi} \intr\int_{\mathbb{R}^{ }} \e^{-\i \omega\kla\left( t \mp r \right)} \, \mathrm{d}^{ }\omega = \frac{\delta (t\mp r)}{4\pi r}. \label{eqn:dAlembert Green solution}
\end{align}
Here we found two qualitatively different Green's functions.
For our purposes the so called \emph{retarded} solution $G=\frac{1}{4\pi r}\delta (t-r)$ will be sufficient -- the other solution is usually referred to as \emph{advanced}.

\section{Exact electromagnetic field of a moving point charge}
Using \autoref{eqn:dAlembert Green solution}, we may now write down the general solution for the four-vector potential 
\begin{align}
A^\mu (\textbf{r},t) &= \frac{\mu_0}{4\pi} \intr\int_{\mathbb{R}^{4}} \frac{J^\mu (\textbf{r}',t') \delta (t - t' - |\textbf{r} - \textbf{r}'|)}{|\textbf{r} - \textbf{r}'|} \, \mathrm{d}^{3}\textbf{r}'\,\tecxt\text{d}t' =  \label{eqn:A mu Greens x4 integral}\\
&= \frac{\mu_0}{4\pi} \intr\int_{\mathbb{R}^{3}} \frac{J^\mu (\textbf{r}', t - |\textbf{r} - \textbf{r}'|)}{|\textbf{r} - \textbf{r}'|} \, \mathrm{d}^{3}\textbf{r}'. \label{eqn:A mu Greens r3 integral}
\end{align}
The second form is often useful, since it already eliminates one of the integrals.
However, for the discussion of a point charge, we will find it easier to work with the form in \autoref{eqn:A mu Greens x4 integral}.
\\
Let $\rho(\textbf{r}, t) = q \delta (\textbf{r} - \textbf{R}(t))$ be the charge density for a given -- but arbitrary -- trajectory $\textbf{R}(t)$.
Then we also have $\textbf{j} = \dot{\textbf{R}} \rho$.
We can compute the electrostatic potential by first integrating over the spatial coordinates
\begin{align}
\phi(\textbf{r},t) &= \frac{\mu_0}{4\pi} \intr\int_{\mathbb{R}^{ }} \intr\int_{\mathbb{R}^{3}} \frac{q\delta (\textbf{r}' - \textbf{R}(t')) \delta (t - t' - |\textbf{r} - \textbf{r}'|)}{|\textbf{r} - \textbf{r}'|} \, \mathrm{d}^{3}\textbf{r}' \, \mathrm{d}^{ }t' = \\
&= \frac{\mu_0 q}{4\pi} \intr\int_{\mathbb{R}^{ }} \frac{\delta \big(f(t')\big)}{|\textbf{r} - \textbf{R}(t')|} \, \mathrm{d}^{ }t'.
\end{align}
Here we defined the scalar function $f(t') := t - t' - |\textbf{r} - \textbf{R}(t')|$.
Note that $f$ is monotonically decreasing, since
\begin{align}
f'(t') &= -1 - \frac{\textbf{r} - \textbf{R}(t')}{|\textbf{r} - \textbf{R}(t')|} \cdot \kla\left(  - \dot{\textbf{R}}(t') \right) \equiv -1 + \textbf{n}\cdot \beta < 0,
\end{align}
where we introduced the unit-vector $\textbf{n} := \frac{\textbf{r} - \textbf{R}(t')}{|\textbf{r} - \textbf{R}(t')|}$ and the velocity $\beta := \dot{\textbf{R}}(t')$ of our particle, for which obviously $|\beta| < 1$.
Also note that $f(t'\rightarrow \pm \infty) \rightarrow \mp \infty$.
Together with the monotony we can conclude that there must exist exactly one $t_r$ such that $f(t_r) = 0$.
With this knowledge, we may now rewrite
\begin{align}
\delta\big(f(t')\big) &= \frac{\delta (t' - t_r)}{|f'(t_r)|} = \frac{\delta (t' - t_r)}{1 - \textbf{n} \cdot \beta}.
\end{align}
Thus the electrostatic potential is given by
\begin{align}
\phi(\textbf{r},t) &= \frac{\mu_0 q}{4\pi} \frac{1}{x - \textbf{x} \cdot \beta}.
\end{align}
Here $\textbf{x} := \textbf{r} - \textbf{R}(t_r)$ and $x := |\textbf{x}|$.
Thus we also have $\textbf{n} = \frac{\textbf{x}}{x}$.
The so called \emph{retarded time} $t_r$ is the solution to the implicit equation $f(t_r) = 0$, which is equivalent to
\begin{align}
t_r &= t - x(t_r).
\end{align}
The vector potential is simply given by $\textbf{A} = \beta \phi$.
\\
To compute the electromagnetic field, we need to compute the derivatives of $t_r$
\begin{align}
\partial_t t_r &= 1 - \partial_t x = 1 - \frac{\textbf{x}}{x} \kla\left( -\beta \right) \partial_t t_r &&\Rightarrow & \partial_tt_r &= \frac{x}{x - \textbf{x}\cdot\beta},\\
\partial_j t_r &= -\partial_j x = -\frac{\textbf{x}}{x} \cdot \kla\left( \textbf{e}_j - \beta \partial_jt_r \right) &&\Rightarrow & \partial_j t_r &= -\frac{\textbf{x}_j}{x - \textbf{x}\cdot \beta}.
\end{align}
From the second equation it also follows that 
\begin{align}
\nabla t_r &= -\frac{\textbf{x}}{x - \textbf{x} \cdot \beta}.
\end{align}
Defining the acceleration $\textbf{a} := \dot{\beta}$ we may now calculate
\begin{align}
-\partial_j \frac{1}{x - \textbf{x}\cdot\beta} &= \frac{1}{\kla\left( x - \textbf{x}\cdot\beta \right)^2} \kla\left( \partial_j x - \partial_j x_k\beta_k  \right) = \frac{-\partial_jt_r - \beta_k \partial_j \textbf{x}_k - \textbf{x}_k \partial_j\beta_k}{(x - \textbf{x}\cdot\beta)^2} = \\
&= \frac{-\partial_j t_r - \beta_k (\delta_{jk} - \beta_k \partial_j t_r) - \textbf{x}_k \textbf{a}_k \partial_j t_r }{(x - \textbf{x}\cdot\beta)^2} = \\
&= \frac{-\beta_j}{(x - \textbf{x}\cdot\beta)^2} + \frac{\textbf{x}_j \kla\left( 1 + \textbf{x}\cdot\textbf{a} - \beta^2 \right)}{(x - \textbf{x}\cdot\beta)^3}. \label{eqn:LW partial j 1 over x}
\end{align}
On the other hand, the time derivative of the same quantity amounts to
\begin{align}
\partial_t \frac{1}{x - \textbf{x} \cdot \beta} &= \frac{\partial_t \textbf{x}_k \beta_k }{\kla\left( x - \textbf{x}\cdot\beta \right)^2} =  \frac{\textbf{x}_k \textbf{a}_k \partial_t t_r + \beta_k (-\beta_k)\partial_tt_r}{\kla\left( x - \textbf{x}\cdot\beta \right)^2}  =  \frac{x \kla\left( \textbf{x}\cdot\textbf{a} - \beta^2 \right)}{\kla\left( x - \textbf{x}\cdot\beta \right)^3}. \label{eqn:LW partial t 1 over x}
\end{align}
Now we have everything ready to compute the electric field
\begin{align}
\textbf{E} &= -\nabla\phi - \partial_t \textbf{A} = \\
&= -\frac{\mu_0 q}{4\pi} \nabla\frac{1}{x - \textbf{x}\cdot \beta} - \beta \frac{\mu_0 q}{4\pi} \partial_t\frac{1}{x - \textbf{x}\cdot \beta} - \frac{\mu_0 q}{4\pi} \frac{1}{x - \textbf{x}\cdot \beta} \textbf{a} \partial_t t_r = \\
&= \frac{\mu_0 q}{4\pi} \klam\left[ \frac{-\beta - x\textbf{a}}{(x - \textbf{x}\cdot \beta)^2} + \frac{\textbf{x}\kla\left( 1 + \textbf{x}\cdot\textbf{a} - \beta^2 \right)}{(x - \textbf{x}\cdot \beta)^3} - \frac{\beta x \kla\left( \textbf{x}\cdot\textbf{a} - \beta^2 \right)}{(x - \textbf{x}\cdot \beta)^3} \right] = \\
&= \frac{\mu_0 q}{4\pi x^2} \klam\left[ -\frac{\beta + x\textbf{a}}{(1 - \textbf{n} \cdot \beta)^2} + \frac{\textbf{n}}{(1 - \textbf{n}\cdot \beta)^3} + \frac{ (\textbf{n} - \beta)\kla\left( \textbf{x}\cdot\textbf{a} - \beta^2 \right)}{(1 - \textbf{n}\cdot \beta)^3} \right].
\end{align}
Note that
\begin{align}
\textbf{n} \times \textbf{E} &= \frac{\mu_0 q}{4\pi} \klam\left[ -\frac{\textbf{n} \times \beta + x\textbf{n} \times \textbf{a}}{(1 - \textbf{n}\cdot \beta)^2} - \frac{\textbf{n}\times \beta \kla\left( \textbf{x}\cdot \textbf{a} - \beta^2 \right)}{(1 - \textbf{n}\cdot \beta)^3} \right] = \\
&= \frac{\mu_0 q}{4\pi} \klam\left[ \frac{\textbf{a}\times \textbf{x}}{(1 - \textbf{n}\cdot \beta)^2} - \frac{\textbf{n}\times\beta \kla\left( 1 - \textbf{n}\cdot\beta + \textbf{x}\cdot \textbf{a} - \beta^2 \right)}{(1 - \textbf{n}\cdot \beta)^3} \right].
\end{align}
Similarly, we may express the magnetic field as
\begin{align}
\textbf{B} &= \nabla\times \textbf{A} = \phi \nabla\times\beta + \nabla\phi \times\beta = \\
&= \phi \kla\left( \nabla t_r \times \textbf{a} \right) + \frac{\mu_0 q}{4\pi } \klam\left[ \frac{\beta}{(x - \textbf{x}\cdot \beta)^2} - \frac{\textbf{x} \kla\left( 1 + \textbf{x}\cdot \textbf{a} - \beta^2 \right)}{(x - \textbf{x}\cdot \beta)^3} \right] \times \beta = \\
&= \frac{\mu_0 q}{4\pi x^2} \klam\left[ \frac{\textbf{a}\times \textbf{x}}{(1 - \textbf{n} \cdot \beta)^2} - \frac{\textbf{n} \times \beta \kla\left( 1 + \textbf{x}\cdot \textbf{a} - \beta^2 \right)}{(1 - \textbf{n}\cdot\beta)^3}  \right] = \\
&= ??? \textbf{n} \times \textbf{E}.
\end{align}

%\begin{thebibliography}{9}
%\bibitem{example} description, \url{https://www.exmapleurl}, day.\,month~2021.
%\end{thebibliography}

\end{document}